Опишите группы автоморфзимов и соответствие между промежуточными полями и подгруппами 

В билете 54 доказывалось, что автоморфизм поля сохраняет $\Q$.

а) $\Q(\sqrt{2})$; 

Минимальным многочленом над $\Q$ для элемента $\sqrt{2}$ является
$x^2 - 2$. Тогда сопряжёнными элементами для $\sqrt{2}$ являются
$\pm \sqrt{2}$. Тогда по предыдущему билету понятно, что будет
два автоморфизма: тождественный и который переводит 
$\sqrt{2}$ в $-\sqrt{2}$.

б) $\Q(\sqrt[3]{2})$; 

Минимальным многочленом над $\Q$ для элемента $\sqrt[3]{2}$ является
$x^3 - 2$. Тогда сопряжённые элементами для $\sqrt[3]{2}$ являются
$\sqrt[3]{2}, \sqrt[3]{2} \omega, \sqrt[3]{2} \omega^2$. 
Тогда по предыдущему билету понятно, что будет один автоморфизм
(тождественный), так как $\sqrt[3]{2} \omega$ и $\sqrt[3]{2} \omega^2$
не лежат в $\Q(\sqrt[3]{2})$.

в) $\Q(\sqrt{2} + \sqrt{3})$.

В билете 40 г) на уд было показано, что минимальным многочленом 
для $\sqrt{2} + \sqrt{3}$ над $\Q$ является
$x^4 - 16x^2 + 1$, а его корнями -- $\pm \sqrt{2} \pm \sqrt{3}$.

По линейности понятно, что $\sqrt{2}$ может отображаться только в
$\pm \sqrt{2}$, а $\sqrt{3}$ -- в $\pm \sqrt{3}$.

Поэтому получаем 4 возможных автоморфизма.
